\begin{explanationbox}
	\textbf{\textcolor{CExplanation}{一句话直觉}}
	当题目中出现 $f'(x)$ 与 $f(x)$ 的不等式时，可以尝试两边除以 $e^x$ 构造 $F(x)=\dfrac{f(x)}{e^x}$，利用其导数符号直接判断单调性。
	\medskip
	\textbf{\textcolor{CExplanation}{核心拆解}}
	\begin{description}[style=nextline, leftmargin=2.8em, labelsep=0.5em, itemsep=0.45em]
		\item[\textbf{要点一（构造动机）：}] 为什么除以 $e^x$？因为 $e^x$ 求导不变，能使导数结果保留 $f'(x)-f(x)$ 的完整结构。
		\item[\textbf{要点二（求导验证）：}] $F(x)=\dfrac{f(x)}{e^x}$，求导得 $F'(x)=\dfrac{f'(x)\cdot e^x - f(x)\cdot e^x}{e^{2x}}$，化简得 $F'(x)=\dfrac{f'(x)-f(x)}{e^x}$。
		\item[\textbf{要点三（符号判单调）：}] 分母 $e^x>0$ 恒成立，所以 $F'(x)$ 的符号完全由 $f'(x)-f(x)$ 决定，从而可直接判断 $F(x)$ 的单调性。
	\end{description}
	\medskip
	\textbf{\textcolor{CExplanation}{几何本质（斜率对应）}}
	函数 $F(x)$ 在某点的导数 $F'(x)$ 表示该点切线的斜率。$F'(x)$ 的正负对应切线倾斜方向，因此 $f'(x)-f(x)$ 的符号决定了 $F(x)$ 图像的上升或下降。
	\medskip
	\textbf{\textcolor{CExplanation}{代数意义（结构转换）}}
	将不易处理的 $f'(x)-f(x)$ 不等式转化为 $F'(x)$ 的符号判断，从而利用单调性比较函数值或求解抽象不等式。
	\medskip
	\textbf{\textcolor{CExplanation}{考点价值（高考必考）}}
	\begin{itemize}[leftmargin=*, itemsep=0.5em, topsep=0.4em]
		\item \textbf{考法一（比较大小）：} 比较 $\dfrac{f(a)}{e^a}$ 与 $\dfrac{f(b)}{e^b}$ 的大小（$a<b$）。
		\item \textbf{考法二（解不等式）：} 已知 $f'(x)-f(x)>0$，解形如 $f(x)>C e^x$ 的不等式，转化为 $F(x)>C$ 的单调性分析。
	\end{itemize}
	\medskip
	\textbf{\textcolor{CExplanation}{顿悟点}}
	为什么选择 $e^x$？因为 $(e^x)'=e^x$，在除法求导后分子恰好是 $f'(x)-f(x)$，没有额外项，从而直接关联条件。
	\medskip
	\textbf{\textcolor{CExplanation}{使用场景}}
	\begin{itemize}[leftmargin=*, itemsep=0.5em, topsep=0.4em]
		\item \textbf{场景一（抽象函数不等式）：} 题干只有 $f(x)$ 的符号信息，没有具体解析式，通过构造单调性解决。
		\item \textbf{场景二（已知导函数关系）：} 已知 $f'(x)$ 与 $f(x)$ 的不等关系，比较 $f(a)e^b$ 与 $f(b)e^a$ 等类似形式。
	\end{itemize}
\end{explanationbox}
